\documentclass[11pt,a4paper]{article}

\usepackage[utf8]{inputenc}
\usepackage[T1]{fontenc}
\usepackage{amsmath,amssymb,amsthm}
\usepackage{mathtools}
\usepackage[margin=1in]{geometry}
\usepackage{enumitem}
\usepackage{xcolor}
\usepackage[colorlinks=true,linkcolor=blue,citecolor=blue]{hyperref}

% Operators
\DeclareMathOperator{\Exp}{Exp}
\DeclareMathOperator{\Log}{Log}
\DeclareMathOperator{\Retr}{Retr}
\DeclareMathOperator{\Image}{Im}
\DeclareMathOperator{\Vol}{Vol}
\DeclareMathOperator{\SE}{SE}
\DeclareMathOperator{\SO}{SO}
\DeclareMathOperator{\so}{\mathfrak{so}}
\DeclareMathOperator{\se}{\mathfrak{se}}
\DeclareMathOperator{\skewop}{skew}
\DeclareMathOperator{\Ad}{Ad}
\DeclareMathOperator{\ad}{ad}

% Commands
\newcommand{\R}{\mathbb{R}}
\newcommand{\E}{\mathbb{E}}
\newcommand{\N}{\mathcal{N}}
\newcommand{\D}{\mathcal{D}}
\newcommand{\Loss}{\mathcal{L}}
\newcommand{\Mphi}{\mathcal{M}_{\phi}}
\newcommand{\Mset}{\mathcal{M}}
\newcommand{\Mpose}{\mathcal{M}_\phi^{\text{pose}}}
\newcommand{\Tphi}{T_\phi}
\newcommand{\Hpose}{H_\phi^{\text{pose}}}
\newcommand{\Jpose}{J_{\text{pose}}}

% Verification flag command
\newcommand{\verifyflag}[1]{\textbf{\color{red}[VERIFY: #1]}}

\title{Core Mathematical Formulation\\
\large Pose-Extended SMCDP: $SE(3)$ End-Effector Self-Model Manifold\\
\normalsize (v2: dimension consistency + sign convention review)}
\author{}
\date{}

\begin{document}

\maketitle

\paragraph{Document version note.} This is v2 of the pose-extended formulation. Changes from v1:
\begin{itemize}[noitemsep]
    \item Replaced ambiguous $(J_H^{\text{pose}})^\top W \, J_H^{\text{pose}}$ with the dimensionally consistent $\Jpose^\top W \Jpose$ formulation, using only the geometric Jacobian $\Jpose$ for metric and DSM computations.
    \item Replaced $\Exp^{-1}_{x_r}(x_0)$ (which would denote the intrinsic Riemannian logarithm on $\Mpose$ under $G_{\text{pose}}$) with the explicitly named ambient/trivialized displacement $\delta x^{\text{amb}}_{r0}$, used as a local approximation to the intrinsic log; the chart target is defined as the weighted least-squares projection of this displacement.
    \item Moved closed-form reward gradients (with $\mathcal{J}_l^{-\top}$ factor) to an appendix and adopted automatic differentiation as the primary implementation, due to sensitivity of sign and Jacobian convention.
    \item Renamed ``ambient dimension'' as ``trivialized tangent / configuration manifold dimension'' to avoid confusion with raw matrix-embedding dimensions of $\SE(3)$.
\end{itemize}

\paragraph{Notation conventions for pose extension.} Throughout this document we use the following primary mathematical spaces:
\begin{itemize}[noitemsep]
    \item $\SE(3) = \SO(3) \ltimes \R^3$: rigid pose configuration space (homogeneous transforms).
    \item $\se(3) \simeq \R^6$: $SE(3)$ Lie algebra, identified with $(\rho, \omega) \in \R^3 \times \R^3$ ($\rho$: linear, $\omega$: angular).
    \item $\so(3) \simeq \R^3$: $SO(3)$ Lie algebra, identified via the hat map $[\,\cdot\,]_\times: \R^3 \to \so(3)$.
    \item All tangent vectors on $SE(3)$ are expressed in \textbf{body-frame (right-trivialized)} coordinates: for $T \in SE(3)$, $T_T \SE(3) \ni \dot T$ is identified with $\xi \in \se(3)$ via $\dot T = T \xi^\wedge$.
\end{itemize}
Quaternions are used only as an implementation representation of $\SO(3)$; all mathematical statements are formulated on $SE(3)$ directly. Conversions and quaternion-specific implementation details are confined to remarks.

\section{Self-model learning}

State: $x = (q, T) \in \R^{n_q} \times \SE(3)$, with $q \in \R^{n_q}$ the joint configuration (chart coordinate) and $T = (R, p) \in \SE(3)$ the end-effector pose ($R \in \SO(3)$, $p \in \R^3$). Embodiment $z_e \in \R^{n_z}$. Goal/condition $c \in \mathcal{C}$ (notation distinct from constraint function $g_\phi$).

\paragraph{Residual forward kinematics on $\SE(3)$.} The analytic forward kinematics produces a pose
\begin{equation}
T_{\text{analytic}}(q, z_e) = \begin{bmatrix} R_{\text{analytic}}(q, z_e) & p_{\text{analytic}}(q, z_e) \\ 0 & 1 \end{bmatrix} \in \SE(3).
\end{equation}
A learned residual twist $\xi_\phi(q, z_e) = (\rho_\phi(q, z_e), \omega_\phi(q, z_e)) \in \se(3)$ is composed via right-multiplication (body-frame residual):
\begin{equation}
\boxed{ \Tphi(q, z_e) = T_{\text{analytic}}(q, z_e) \cdot \exp_{\SE(3)}\!\bigl(\xi_\phi(q, z_e)^\wedge\bigr) }
\end{equation}
We denote the induced position and rotation components by $p_\phi(q, z_e), R_\phi(q, z_e)$, so that $\Tphi(q, z_e) = (R_\phi, p_\phi)$.

\paragraph{Frame convention.} The right-multiplicative form models the residual as a body-frame perturbation of the analytic end-effector pose, which is the natural model for tool/end-effector calibration drift, gripper compliance, and tool-mounting error. Throughout this document, $\xi_\phi$ is in body-frame; all subsequent uses of $\log_{\SE(3)}$, $\log_{\SO(3)}$, $J_\omega$ are accordingly body-frame-trivialized.

\paragraph{Constraint function on $SE(3)$.} Define
\begin{equation}
g_\phi(x, z_e) = \Log_{\SE(3)}\!\Bigl( \Tphi(q, z_e)^{-1} T \Bigr) \in \se(3),
\end{equation}
where $\Log_{\SE(3)}: \SE(3) \to \se(3)$ is the principal logarithm. Equivalently, splitting into translation/rotation components in body-frame:
\begin{equation}
g_\phi(x, z_e) = \begin{bmatrix} R_\phi(q, z_e)^\top \bigl(p - p_\phi(q, z_e)\bigr) \\ \Log_{\SO(3)}\!\bigl( R_\phi(q, z_e)^\top R \bigr) \end{bmatrix}.
\end{equation}
Note $g_\phi(x, z_e) = 0 \iff T = \Tphi(q, z_e)$.

\paragraph{Self-model loss.} Let $\D_{\text{self}}$ be the self-exploration dataset of $(q, T_{\text{true}}, z_e)$ tuples. Define the body-frame error twist
\begin{equation}
e_{\text{self}}(q, T_{\text{true}}, z_e) = \Log_{\SE(3)}\!\Bigl( \Tphi(q, z_e)^{-1} T_{\text{true}} \Bigr) = (e_\rho, e_\omega) \in \R^3 \times \R^3.
\end{equation}
The self-model loss is
\begin{equation}
\Loss_{\text{self}}(\phi) = \E_{(q, T_{\text{true}}, z_e) \sim \D_{\text{self}}} \bigl[\, w_p \| e_\rho \|^2 + w_R \| e_\omega \|^2 \,\bigr] + \beta_p \, \E\bigl[\| \nabla_q \rho_\phi \|_F^2\bigr] + \beta_R \, \E\bigl[\| \nabla_q \omega_\phi \|_F^2\bigr]
\end{equation}
The fidelity weights $w_p, w_R > 0$ reconcile the unit mismatch (meters vs.\ radians); the smoothness weights $\beta_p, \beta_R > 0$ regularize $\rho_\phi, \omega_\phi$ separately.

\paragraph{Position-only specialization.} Setting $\omega_\phi \equiv 0$ and $w_R \to 0$ reduces the self-model to translation-only. The induced position component is
\begin{equation}
p_\phi(q, z_e) = p_{\text{analytic}}(q, z_e) + R_{\text{analytic}}(q, z_e) \, \rho_\phi(q, z_e),
\end{equation}
so $\rho_\phi$ is a body-frame translational residual. Recovering the original position-only framework with world-frame residual $\Delta_\phi$ requires the change of variables $\Delta_\phi(q, z_e) = R_{\text{analytic}}(q, z_e) \, \rho_\phi(q, z_e)$. The pose-extended framework thus generalizes the position-only one up to the choice of residual frame; the two are related by a $q, z_e$-dependent rotation.

\section{Manifold from self-model}

\paragraph{Self-model graph manifold.}
\begin{equation}
\Mpose(z_e) = \bigl\{ (q, T) \in \R^{n_q} \times \SE(3) \;:\; T = \Tphi(q, z_e) \bigr\}.
\end{equation}
This is an $n_q$-dimensional submanifold of $\R^{n_q} \times \SE(3)$, parameterized globally by the chart coordinate $q$.

\paragraph{Embedding map.}
\begin{equation}
\Hpose(q, z_e) = \bigl(q, \, \Tphi(q, z_e)\bigr) \in \R^{n_q} \times \SE(3).
\end{equation}

\paragraph{Tool-tip instantiation.} For a manipulator with embodiment $z_e \in \R$ representing tool length along the end-effector $z$-axis,
\begin{equation}
T_{\text{analytic}}(q, z_e) = T_{\text{hand}}(q) \cdot \begin{bmatrix} I_3 & (0, 0, z_e)^\top \\ 0 & 1 \end{bmatrix},
\end{equation}
where $T_{\text{hand}}(q) = (R_{\text{hand}}(q), p_{\text{hand}}(q)) \in \SE(3)$ is the hand-link pose. The residual $\xi_\phi(q, z_e)$ captures compliance, calibration drift, tool-mounting error, and other analytic-uncapturable effects.

\section{Differential structure (tangent bundle)}
\label{sec:tangent}

\paragraph{Geometric Jacobian (body-frame).} For $\Tphi(q, z_e) \in \SE(3)$, the body-frame velocity twist induced by chart velocity $\dot q \in \R^{n_q}$ is
\begin{equation}
\xi_{\text{body}}(q, \dot q, z_e) = \Jpose(q, z_e) \, \dot q \in \se(3),
\qquad
\Jpose(q, z_e) = \begin{bmatrix} J_v(q, z_e) \\ J_\omega(q, z_e) \end{bmatrix} \in \R^{6 \times n_q},
\end{equation}
with body-frame linear and angular Jacobians $J_v, J_\omega \in \R^{3 \times n_q}$, defined by $\dot \Tphi = \Tphi (\Jpose \dot q)^\wedge$.

\paragraph{Tangent map (lift) at a manifold point.} For $x = \Hpose(q, z_e) \in \Mpose(z_e)$, the embedding differential is the linear map
\begin{equation}
J_H^{\text{pose}}(q, z_e) : \R^{n_q} \longrightarrow \R^{n_q} \times \se(3), \qquad \dot q \;\longmapsto\; \bigl(\dot q, \, \Jpose(q, z_e) \dot q\bigr).
\end{equation}
In matrix form (with codomain identified as $\R^{n_q} \times \R^6$ via body-frame trivialization):
\begin{equation}
J_H^{\text{pose}}(q, z_e) = \begin{bmatrix} I_{n_q} \\ \Jpose(q, z_e) \end{bmatrix} \in \R^{(n_q + 6) \times n_q}.
\end{equation}

\paragraph{Important interpretation.} $J_H^{\text{pose}}$ is a tangent map into the trivialized tangent space $\R^{n_q} \times \se(3)$, \emph{not} a Jacobian into raw quaternion or rotation matrix coordinates. Body-frame trivialization identifies $T_T \SE(3)$ with $\R^6$ via $\dot T = T \xi^\wedge$, so the codomain of $J_H^{\text{pose}}$ has trivialized tangent dimension $n_q + 6$, distinct from any matrix-embedding dimension of $\SE(3)$ itself ($\R^{12}$ for $(R, p)$, $\R^7$ for quaternion implementation, etc.).

\paragraph{Tangent space at a manifold point.}
\begin{equation}
T_x \Mpose = \Image J_H^{\text{pose}}(q, z_e) = \bigl\{ (\dot q, \, \Jpose \dot q) : \dot q \in \R^{n_q} \bigr\} \;\subset\; \R^{n_q} \times \se(3).
\end{equation}

\paragraph{Constraint Jacobian and on-manifold tangent verification.} The constraint $g_\phi(x, z_e) = \Log_{\SE(3)}(\Tphi(q, z_e)^{-1} T)$ has differential at points $x \in \Mpose(z_e)$ (i.e., where $T = \Tphi(q, z_e)$), in body-frame trivialization:
\begin{equation}
J_g(x, z_e) \big|_{\Mpose} = \begin{bmatrix} -\Jpose(q, z_e) & I_6 \end{bmatrix} \in \R^{6 \times (n_q + 6)},
\end{equation}
where the column blocks correspond to $(\delta q, \xi) \in \R^{n_q} \times \se(3)$.

\paragraph{On-manifold tangent identity.}
\begin{equation}
\boxed{ J_g(x, z_e) \cdot J_H^{\text{pose}}(q, z_e) = -\Jpose + \Jpose = 0 \quad \text{on } \Mpose(z_e). }
\end{equation}
Thus any vector of the form $J_H^{\text{pose}} a$ for $a \in \R^{n_q}$ lies in the null space of the on-manifold linearization of $g_\phi$, i.e., in $T_x \Mpose$.

\paragraph{Off-manifold remark.} The full nonlinear constraint $g_\phi$ involves $\Log_{\SE(3)}$, whose Jacobian off-manifold contains additional right-Jacobian factors ($\mathcal{J}_l^{-1}(\xi)$ for $\xi \neq 0$). These factors evaluate to the identity at $\xi = 0$, recovering the on-manifold linearization above. Since the framework's sampling generates $x$ via the embedding $\Hpose$ at every step (Section~\ref{sec:sampling}), all tangent verifications used in practice are evaluated on-manifold.

\section{Induced Riemannian metric}

\paragraph{Pose-tangent inner product (on $\se(3)$).} The Lie algebra $\se(3)$ carries the weighted inner product
\begin{equation}
\langle \xi_1, \xi_2 \rangle_W = \xi_1^\top W \xi_2, \qquad W = \begin{bmatrix} W_p & 0 \\ 0 & W_R \end{bmatrix} \in \R^{6 \times 6},
\end{equation}
with positive-definite weights $W_p, W_R \in \R^{3 \times 3}$ reconciling translation/rotation units. Standard choice: $W_p = \sigma_p^{-2} I_3$, $W_R = \sigma_R^{-2} I_3$.

\paragraph{Trivialized ambient inner product.} On $\R^{n_q} \times \se(3)$:
\begin{equation}
\bigl\langle (\dot q_1, \xi_1), \, (\dot q_2, \xi_2) \bigr\rangle = \dot q_1^\top \dot q_2 + \xi_1^\top W \xi_2.
\end{equation}
Note that the chart-coordinate part uses the standard Euclidean inner product on $\R^{n_q}$ (no weighting), while the pose-tangent part uses $W$.

\paragraph{Chart-coordinate metric (induced).} For chart vectors $a, b \in \R^{n_q}$ lifted as $u = J_H^{\text{pose}} a = (a, \Jpose a)$, $v = J_H^{\text{pose}} b = (b, \Jpose b)$:
\begin{equation}
\bigl\langle u, v \bigr\rangle_{\Mset, x} = a^\top b + (\Jpose a)^\top W (\Jpose b) = a^\top \bigl( I_{n_q} + \Jpose^\top W \Jpose \bigr) b.
\end{equation}
Therefore (\textbf{dimensionally consistent form, replacing v1's $(J_H^{\text{pose}})^\top W J_H^{\text{pose}}$}):
\begin{equation}
\boxed{ G_{\text{pose}}(q, z_e) = I_{n_q} + \Jpose(q, z_e)^\top W \, \Jpose(q, z_e) = I_{n_q} + J_v^\top W_p J_v + J_\omega^\top W_R J_\omega. }
\end{equation}
Here $\Jpose \in \R^{6 \times n_q}$ and $W \in \R^{6 \times 6}$, so the product $\Jpose^\top W \Jpose \in \R^{n_q \times n_q}$ is well-defined.

\paragraph{Norm equivalence.}
\begin{equation}
\bigl\| (a, \Jpose a) \bigr\|^2 = a^\top G_{\text{pose}}(q, z_e) \, a = \| a \|_{G_{\text{pose}}}^2.
\end{equation}

\paragraph{Riemannian gradient (chart).} For $f : \Mpose \to \R$ with chart pullback $\bar f(q) = f(\Hpose(q, z_e))$:
\begin{equation}
\nabla_{\Mset} f \;\leftrightarrow\; G_{\text{pose}}(q, z_e)^{-1} \nabla_q \bar f.
\end{equation}

\paragraph{Density convention.} All densities $p_r$ on $\Mpose(z_e)$ are defined with respect to the Riemannian volume measure $d\Vol_{\Mset} = \sqrt{\det G_{\text{pose}}(q, z_e)} \, dq$.

\paragraph{Position-only consistency.} Setting $W_R = 0$ reduces $G_{\text{pose}} = I + J_v^\top W_p J_v$, recovering the position-only metric (with $W_p = I$ and $J_v = J_F$).

\section{Trajectory product manifold}
\label{sec:product-manifold}

Trajectory $\tau = (x_0, x_1, \ldots, x_H)$, $x_h \in \Mpose(z_e)$:
\begin{equation}
\tau \in \Mpose(z_e)^{H+1}.
\end{equation}
The trajectory chart is $q_{0:H} = (q_0, \ldots, q_H) \in \R^{n_q (H+1)}$; the manifold dimension is $n_q(H+1)$, while the trivialized tangent / configuration manifold dimension grows to $(n_q + 6)(H+1)$.

\paragraph{Product tangent space and metric.}
\begin{equation}
T_\tau \Mpose^{H+1} = \prod_{h=0}^{H} T_{x_h} \Mpose, \qquad \langle U, V \rangle_\tau = \sum_{h=0}^{H} \langle u_h, v_h \rangle_{\Mset, x_h}.
\end{equation}
The product chart metric is block-diagonal:
\begin{equation}
G_{\text{traj}}^{\text{pose}}(q_{0:H}, z_e) = \operatorname{diag}\bigl( G_{\text{pose}}(q_0, z_e), \ldots, G_{\text{pose}}(q_H, z_e) \bigr) \in \R^{n_q(H+1) \times n_q(H+1)}.
\end{equation}

\section{Riemannian SDE}

Diffusion time $r \in [0, 1]$ (distinct from trajectory index $h$).

\paragraph{Forward SDE on $\Mpose$.}
\begin{equation}
dX_r = b(X_r) \, dr + dB^{\Mset}_r, \qquad b(X_r) = -\tfrac{1}{2} \nabla_{\Mset} U(X_r),
\end{equation}
where $B^{\Mset}_r$ is Brownian motion on $\Mpose(z_e)$ with respect to the Riemannian volume measure.

\paragraph{Langevin potential with Riemannian volume correction.}
\begin{equation}
\boxed{ U^{\text{pose}}(q, z_e) = \frac{1}{2 \gamma^2} \| q - \mu_q \|^2 + \frac{1}{2} \log \det G_{\text{pose}}(q, z_e). }
\end{equation}

\paragraph{Chart-form drift.}
\begin{equation}
b^q(q, z_e) = -\tfrac{1}{2} G_{\text{pose}}^{-1} \nabla_q U^{\text{pose}} = -\frac{1}{2 \gamma^2} G_{\text{pose}}^{-1}(q - \mu_q) - \frac{1}{4} G_{\text{pose}}^{-1} \nabla_q \log \det G_{\text{pose}}.
\end{equation}

\paragraph{Reverse SDE.}
\begin{equation}
dY_r = \bigl[ -b(Y_r) + \nabla_{\Mset} \log p_{1-r}(Y_r) \bigr] dr + dB^{\Mset}_r.
\end{equation}

\paragraph{Trajectory-level on product manifold.}
\begin{equation}
d\tau_r = b(\tau_r) \, dr + dB^{\Mset^{H+1}}_r.
\end{equation}

\section{Score model}

\paragraph{Score field as tangent bundle section.}
\begin{equation}
s_\theta(r, x, c, z_e) \in T_x \Mpose.
\end{equation}

\paragraph{Chart-coordinate output with self-model lift.} The score network outputs a chart vector $s_\theta^q \in \R^{n_q}$, and the trivialized ambient score is obtained by lift:
\begin{equation}
s_\theta^{\text{amb}}(r, x, c, z_e) = J_H^{\text{pose}}(q, z_e) \, s_\theta^q(r, q, c, z_e) = \bigl(s_\theta^q, \, \Jpose(q, z_e) \, s_\theta^q\bigr) \in \R^{n_q} \times \se(3).
\end{equation}

\paragraph{Tangent verification.} By the on-manifold identity $J_g \, J_H^{\text{pose}} = 0$:
\begin{equation}
J_g(x, z_e) \, s_\theta^{\text{amb}} = J_g \, J_H^{\text{pose}} \, s_\theta^q = 0 \quad \text{on } \Mpose(z_e).
\end{equation}

\section{Score matching loss (Riemannian DSM)}
\label{sec:dsm}

Varadhan asymptotic (small $r$):
\begin{equation}
\nabla_{\Mset} \log p_{r \mid 0}(x_r \mid x_0) \;\approx\; \frac{1}{r} \, \mathrm{Log}_{\Mset, x_r}(x_0),
\end{equation}
where $\mathrm{Log}_{\Mset, x_r}: \Mpose \to T_{x_r} \Mpose$ is the intrinsic Riemannian logarithm on the graph manifold $\Mpose$ under metric $G_{\text{pose}}$.

\paragraph{Local approximation: trivialized ambient displacement.} The intrinsic log $\mathrm{Log}_{\Mset, x_r}(x_0)$ has no closed form on $\Mpose$ in general. Following the same pattern as the position-only formulation, we use the local approximation
\begin{equation}
\boxed{ \delta x^{\text{amb}}_{r0} = \begin{pmatrix} q_0 - q_r \\ \Log_{\SE(3)}\!\bigl( \Tphi(q_r, z_e)^{-1} \Tphi(q_0, z_e) \bigr) \end{pmatrix} \in \R^{n_q} \times \se(3) }
\end{equation}
as a stand-in for the intrinsic log. The first component is the chart-space (Euclidean) displacement; the second component is the body-frame $\SE(3)$ logarithm. Note that $\delta x^{\text{amb}}_{r0}$ is generally \emph{not} an element of $T_{x_r} \Mpose = \Image J_H^{\text{pose}}$.

\paragraph{Validity.} The approximation $\delta x^{\text{amb}}_{r0} \approx \mathrm{Log}_{\Mset, x_r}(x_0)$ holds:
\begin{itemize}[noitemsep]
    \item In the local regime where $q_r \approx q_0$ (small DSM step, small $r$);
    \item When $G_{\text{pose}}$ is slowly varying (so the metric is approximately constant over the displacement);
    \item When the $\SE(3)$ component itself is small ($\Log_{\SE(3)}$ argument near identity).
\end{itemize}
This is the same approximation used (implicitly) in the position-only framework. We make it explicit here to clarify what $\Exp^{-1}$ denotes operationally.

\paragraph{Chart target via weighted least-squares projection.} The chart-coordinate target $a^*_{\text{pose}}(q_r, x_0) \in \R^{n_q}$ is defined as the $\widetilde W$-weighted least-squares projection of $\delta x^{\text{amb}}_{r0}$ onto $\Image J_H^{\text{pose}}$:
\begin{equation}
a^*_{\text{pose}}(q_r, x_0) = \arg\min_{a \in \R^{n_q}} \bigl\| J_H^{\text{pose}}(q_r, z_e) \, a - \delta x^{\text{amb}}_{r0} \bigr\|_{\widetilde W}^2,
\end{equation}
where $\widetilde W = \mathrm{diag}(I_{n_q}, W) \in \R^{(n_q + 6) \times (n_q + 6)}$.

\paragraph{Closed form.} Expanding the normal equations using $(J_H^{\text{pose}})^\top \widetilde W J_H^{\text{pose}} = G_{\text{pose}}$:
\begin{equation}
\boxed{ a^*_{\text{pose}}(q_r, x_0) = G_{\text{pose}}(q_r, z_e)^{-1} \bigl[ (q_0 - q_r) + \Jpose(q_r, z_e)^\top W \, \Log_{\SE(3)}\!\bigl( \Tphi(q_r)^{-1} \Tphi(q_0) \bigr) \bigr]. }
\end{equation}
This block-decomposed form makes explicit how the chart and pose components contribute separately to the chart target.

\paragraph{Important relation.} We have
\begin{equation}
J_H^{\text{pose}} a^*_{\text{pose}} = \Pi_{\Image J_H^{\text{pose}}}^{\widetilde W} \bigl( \delta x^{\text{amb}}_{r0} \bigr),
\end{equation}
i.e., $J_H^{\text{pose}} a^*_{\text{pose}}$ is the orthogonal projection (under $\widetilde W$) of $\delta x^{\text{amb}}_{r0}$ onto the on-manifold tangent image. The equality $J_H^{\text{pose}} a^*_{\text{pose}} = \delta x^{\text{amb}}_{r0}$ does \emph{not} hold in general, since $\delta x^{\text{amb}}_{r0} \notin \Image J_H^{\text{pose}}$.

Three forms of the loss:

\textbf{(a) Ambient (full Riemannian, exact under intrinsic-log assumption):}
\begin{equation}
\Loss^{\text{ambient}} = \E_{r, x_0, x_r} \left[ \left\| s_\theta(r, x_r, c, z_e) - \frac{1}{r} \mathrm{Log}_{\Mset, x_r}(x_0) \right\|^2_{x_r, W} \right]
\end{equation}
where $\| \cdot \|_{x_r, W}$ is the $W$-weighted ambient pose-tangent norm.

\textbf{(b) Chart $G$-weighted (used in practice, with $\delta x^{\text{amb}}$ approximation):}
\begin{equation}
\Loss^{\text{chart-}G} = \E_{r, q_0, q_r} \left[ (s_\theta^q - a^*_{\text{pose}})^\top G_{\text{pose}}(q_r, z_e) \, (s_\theta^q - a^*_{\text{pose}}) \right].
\end{equation}

Equivalence (under $\delta x^{\text{amb}} \approx \mathrm{Log}_{\Mset, x_r}$):
\begin{equation}
\| J_H^{\text{pose}} s_\theta^q - J_H^{\text{pose}} a^*_{\text{pose}} \|^2_{\widetilde W} = (s_\theta^q - a^*_{\text{pose}})^\top G_{\text{pose}} \, (s_\theta^q - a^*_{\text{pose}}).
\end{equation}

\textbf{(c) Chart-Euclidean (computational approximation, not exact Riemannian DSM):}
\begin{equation}
\Loss^{\text{chart-Eucl}} \approx \E_{r, q_0, q_r} \left[ \left\| s_\theta^q - \frac{q_0 - q_r}{r} \right\|^2 \right].
\end{equation}
Valid when $G_{\text{pose}} \approx I$ or slowly varying, and the pose displacement is small.

\section{Sampling: retraction-based geodesic random walk}
\label{sec:sampling}

\paragraph{Implementation choice.} The framework uses retraction-based geodesic random walk (manifold-direct simulation) as the main sampling method. This avoids explicit chart-coordinate It\^o SDE simulation.

\paragraph{Graph retraction (pose).}
\begin{equation}
\Retr_{(q, T)}(\delta q) = \bigl( q + \delta q, \, \Tphi(q + \delta q, z_e) \bigr).
\end{equation}
The pose component is recomputed from the updated chart $q + \delta q$ via the self-model. Since $\Tphi(\cdot, z_e) \in \SE(3)$ by construction, the result lies on $\Mpose(z_e)$ exactly.

\paragraph{Tangent noise (Riemannian Gaussian).}
\begin{equation}
\xi_k^q \sim \N\bigl(0, \, G_{\text{pose}}(q_k, z_e)^{-1}\bigr).
\end{equation}

\paragraph{Reverse sampling step (component-wise on product manifold).}
\begin{align}
\mu_{h, k}^q &= -b^q(q_{h, k}) + s_{\theta, h}^q(r_k, \tau_k, c, z_e) \\
\delta q_{h, k} &= \Delta r \cdot \mu_{h, k}^q + \sqrt{\Delta r} \cdot \xi_{h, k}^q \\
q_{h, k+1} &= q_{h, k} + \delta q_{h, k} \\
x_{h, k+1} &= \Hpose(q_{h, k+1}, z_e).
\end{align}
Manifold adherence holds at every step by construction:
\begin{equation}
\| g_\phi(x_{h, k+1}, z_e) \| = 0 \quad \text{(by construction, exact)}.
\end{equation}

\paragraph{Reference distribution.}
\begin{equation}
p_{\text{ref}}(\tau \mid z_e) \propto \exp\!\left( -\sum_{h=0}^{H} \frac{d_{\Mset}(x_h, \mu)^2}{2 \gamma^2} \right),
\end{equation}
approximated by chart-coordinate Gaussian $q_h \sim \N(\mu_q, \gamma^2 G_{\text{pose}}(\mu_q, z_e)^{-1})$, then lifting via $\Hpose$.

\paragraph{Implementation remark on rotation representation.} While $\Tphi(q, z_e) \in \SE(3)$ mathematically, implementations typically store $R_\phi$ as a unit quaternion. The unit-norm constraint is preserved exactly by composing through $\Tphi$ at every step rather than by additive update on the quaternion. Double-cover ambiguity is irrelevant since rotation operations go through $R$ or $\Log_{\SO(3)}$. For trajectory storage in quaternion form $(\bar q_{R, h})_{h=0}^H$, we additionally enforce \emph{hemisphere consistency} by requiring $\bar q_{R, h}^\top \bar q_{R, h-1} \geq 0$ along the trajectory (flipping $\bar q_{R, h} \to -\bar q_{R, h}$ otherwise); this prevents arbitrary sign jumps in stored quaternion sequences without affecting any rotation operation.

\section{Goal-conditioned guidance (single-point, framework template)}

Conditional score:
\begin{equation}
s_\theta(r, x, c, z_e) \approx \nabla_{\Mset} \log p_r(x \mid c, z_e).
\end{equation}

\paragraph{Reward-based guidance (single-point, chart form).}
\begin{equation}
s_{\text{guided}}^q = s_\theta^q + \alpha \, G_{\text{pose}}(q, z_e)^{-1} \nabla_q \bar R(q, c).
\end{equation}

\paragraph{Implementation note.} In practice, $\nabla_q \bar R(q, c)$ is computed via automatic differentiation through the reward function (which itself depends on $\Tphi(q, z_e)$ and the target). This avoids manual derivation of the right-Jacobian factors of $\Log_{\SE(3)}$ and the associated convention issues. Closed-form expressions are documented in Appendix~\ref{app:closed-form-grad} for reference and verification, but the main implementation uses autograd.

\section{Application-specific extension (trajectory-level guidance)}

\subsection{Condition extension}

For pose-conditional manipulation tasks:
\begin{equation}
c = (T_{\text{start}}, T_{\text{target}}, z_e), \qquad T_{\text{start}}, T_{\text{target}} \in \SE(3).
\end{equation}
Implementation: $T_\bullet$ parameterized as $(p_\bullet, \bar q_{R, \bullet}) \in \R^3 \times S^3$ for network input.

\subsection{Trajectory-level Riemannian gradient}

For trajectory $\tau \in \Mpose^{H+1}$ and reward function $R: \Mpose^{H+1} \times \mathcal{C} \to \R$:
\begin{equation}
\nabla_{\Mset^{H+1}} R(\tau) \;\leftrightarrow\; G_{\text{traj}}^{\text{pose}}(q_{0:H}, z_e)^{-1} \nabla_{q_{0:H}} \bar R(\tau).
\end{equation}

\subsection{Trajectory-level guidance}

\begin{equation}
\boxed{ s_{\text{guided}}^q(r, \tau, c, z_e) = s_\theta^q(r, \tau, c, z_e) + G_{\text{traj}}^{\text{pose}}(q_{0:H}, z_e)^{-1} \nabla_{q_{0:H}} R_{\text{total}}(\tau, c). }
\end{equation}
The gradient $\nabla_{q_{0:H}} R_{\text{total}}$ is computed via automatic differentiation; closed-form expressions are in Appendix~\ref{app:closed-form-grad}.

\subsection{Multi-component pose reward}

\begin{equation}
R_{\text{total}}(\tau, c) = \alpha_s R_{\text{start}}(\tau, c) + \alpha_g R_{\text{goal}}(\tau, c) + \alpha_v R_{\text{vel}}(\tau) + \alpha_a R_{\text{acc}}(\tau).
\end{equation}

\paragraph{Pose start anchoring.}
\begin{equation}
e_{\text{start}}(q_0, z_e) = \Log_{\SE(3)}\!\bigl( \Tphi(q_0, z_e)^{-1} T_{\text{start}} \bigr) = (e_\rho^{\text{s}}, e_\omega^{\text{s}}),
\end{equation}
\begin{equation}
R_{\text{start}}(\tau, c) = -\bigl( \alpha_p^{\text{s}} \| e_\rho^{\text{s}} \|^2 + \alpha_R^{\text{s}} \| e_\omega^{\text{s}} \|^2 \bigr).
\end{equation}

\paragraph{Pose endpoint guidance.}
\begin{equation}
e_{\text{goal}}(q_H, z_e) = \Log_{\SE(3)}\!\bigl( \Tphi(q_H, z_e)^{-1} T_{\text{target}} \bigr) = (e_\rho^{\text{g}}, e_\omega^{\text{g}}),
\end{equation}
\begin{equation}
R_{\text{goal}}(\tau, c) = -\bigl( \alpha_p^{\text{g}} \| e_\rho^{\text{g}} \|^2 + \alpha_R^{\text{g}} \| e_\omega^{\text{g}} \|^2 \bigr).
\end{equation}

\paragraph{Decoupled position/rotation reward (alternative).}
\begin{equation}
R_{\text{goal}}^{\text{decoupled}} = -\alpha_p \| p_\phi(q_H, z_e) - p_{\text{target}} \|^2 - \alpha_R \| \Log_{\SO(3)}\!\bigl( R_\phi(q_H, z_e)^\top R_{\text{target}} \bigr) \|^2.
\end{equation}
Note: position term in world-frame ($\R^3$ Euclidean), rotation term in body-frame trivialized $\SO(3)$ log. Either decoupled or coupled twist-error form is admissible.

\paragraph{Velocity smoothness (chart-space).}
\begin{equation}
R_{\text{vel}}(\tau) = -\sum_{h=0}^{H-1} \| q_{h+1} - q_h \|^2.
\end{equation}

\paragraph{Acceleration smoothness.}
\begin{equation}
R_{\text{acc}}(\tau) = -\sum_{h=1}^{H-1} \| q_{h+1} - 2 q_h + q_{h-1} \|^2.
\end{equation}

\subsection{Tangent verification (trajectory-level)}

For each $h$, on $\Mpose$:
\begin{equation}
J_g(x_h, z_e) \cdot \bigl[ J_H^{\text{pose}}(q_h, z_e) \, s_{\text{guided}, h}^q \bigr] = 0,
\end{equation}
so the lifted trajectory score lies in $T_\tau \Mpose^{H+1}$.

\subsection{Timestep-selective guidance application}

Mask $m_h \in \{0, 1\}$:
\begin{equation}
s_{\text{guided}, h}^q = s_{\theta, h}^q + m_h \cdot \bigl[ G_{\text{pose}}(q_h, z_e)^{-1} \nabla_{q_h} R_{\text{total}}(\tau, c) \bigr].
\end{equation}

Common mask schedules:
\begin{align}
m_h^{\text{all}} &= 1 \quad \forall h, &
m_h^{\text{last\_only}} &= \mathbb{1}[h = H], \\
m_h^{\text{last\_half}} &= \mathbb{1}[h \geq H/2], &
m_h^{\text{last\_quarter}} &= \mathbb{1}[h \geq 3H/4].
\end{align}

\subsection{Reverse sampling step with guidance}

\begin{align}
\mu_{h, k}^q &= -b^q(q_{h, k}) + s_{\theta, h}^q(r_k, \tau_k, c, z_e) + m_h \cdot G_{\text{pose}}(q_{h, k}, z_e)^{-1} \nabla_{q_h} R_{\text{total}}(\tau_k, c) \\
\delta q_{h, k} &= \Delta r \cdot \mu_{h, k}^q + \sqrt{\Delta r} \cdot \xi_{h, k}^q, \qquad \xi_{h, k}^q \sim \N(0, G_{\text{pose}}(q_{h, k}, z_e)^{-1}) \\
q_{h, k+1} &= q_{h, k} + \delta q_{h, k} \\
x_{h, k+1} &= \Hpose(q_{h, k+1}, z_e).
\end{align}

\section{Evaluation metrics}

\paragraph{Manifold adherence (by construction).}
\begin{equation}
\max_h \| g_\phi(x_h, z_e) \|_W = 0.
\end{equation}

\paragraph{Endpoint position error.}
\begin{equation}
e_p(\tau) = \| p_\phi(q_H, z_e) - p_{\text{target}} \|.
\end{equation}

\paragraph{Endpoint rotation error (geodesic).}
\begin{equation}
e_R(\tau) = \bigl\| \Log_{\SO(3)}\!\bigl( R_\phi(q_H, z_e)^\top R_{\text{target}} \bigr) \bigr\| \in [0, \pi].
\end{equation}

\paragraph{Pose success rate at $(\rho_p, \rho_R)$.}
\begin{equation}
\text{succ}@(\rho_p, \rho_R) = \E_{\tau \sim p_\theta} \bigl[ \mathbb{1}\{ e_p(\tau) \leq \rho_p \;\land\; e_R(\tau) \leq \rho_R \} \bigr].
\end{equation}

\paragraph{Trajectory smoothness (joint-space).}
\begin{equation}
\text{vel}(\tau) = \frac{1}{H} \sum_{h=0}^{H-1} \| q_{h+1} - q_h \|, \qquad \text{acc}(\tau) = \frac{1}{H-1} \sum_{h=1}^{H-1} \| q_{h+1} - 2 q_h + q_{h-1} \|.
\end{equation}

\paragraph{Mode capture, Wasserstein, joint violation, quaternion norm sanity.} Same as position-only formulation; omitted for brevity.

\section{Pipeline summary}

\subsection{Framework pipeline (geometric core, pose-extended)}

\begin{align}
\text{Stage 1.}\quad & \Tphi = T_{\text{analytic}} \cdot \exp_{\SE(3)}(\xi_\phi^\wedge) \quad \text{via } \min_\phi \Loss_{\text{self}} \\
\text{Stage 2.}\quad & \Mpose(z_e) = \{(q, T) : T = \Tphi(q, z_e)\}, \quad G_{\text{pose}} = I + \Jpose^\top W \Jpose \\
\text{Stage 3.}\quad & \tau_i = (\Hpose(q_h, z_e))_{h=0}^H \in \Mpose^{H+1} \\
\text{Stage 4.}\quad & d\tau_r = b(\tau_r) \, dr + dB^{\Mset^{H+1}}_r \quad \text{(with } \tfrac{1}{2}\log\det G_{\text{pose}} \text{ correction)} \\
\text{Stage 5.}\quad & a^*_{\text{pose}} = G_{\text{pose}}^{-1} \bigl[ (q_0 - q_r) + \Jpose^\top W \Log_{\SE(3)}(\Tphi(q_r)^{-1} \Tphi(q_0)) \bigr] \\
\text{Stage 6.}\quad & \tau_K \sim p_{\text{ref}} \;\to\; \tau_0 \in \Mpose(z_e)^{H+1}.
\end{align}

\subsection{Application-specific extension}

\begin{align}
\text{Stage 5'.}\quad & c = (T_{\text{start}}, T_{\text{target}}, z_e), \quad \text{channel-concat} \\
\text{Stage 6'.}\quad & s_{\text{guided}}^q = s_\theta^q + (G_{\text{traj}}^{\text{pose}})^{-1} \nabla_{q_{0:H}} R_{\text{total}}(\tau, c) \quad \text{(via autograd)}.
\end{align}

\section{Future work: physically-grounded weight matrix \texorpdfstring{$W$}{W}}
\label{sec:future-W}

In the present formulation, $W = \mathrm{diag}(W_p, W_R)$ is introduced as a unit-reconciling diagonal weight on $\se(3)$, with $W_p = \sigma_p^{-2} I_3$ (m$^{-2}$) and $W_R = \sigma_R^{-2} I_3$ (rad$^{-2}$). This choice is sufficient for dimensional consistency but treats $W$ as a free hyperparameter rather than a derived object. The mathematical structure of the framework, however, accommodates any positive-definite $W$ (possibly $q$-dependent) on $\se(3)$, opening natural connections to robot-physics-grounded choices that extend beyond unit reconciliation. We outline three such directions and summarize their implications.

\subsection{Choice (i): manipulability tensor}

The body-frame end-effector twist Jacobian $\Jpose(q, z_e)$ defines the velocity manipulability ellipsoid (Yoshikawa, 1985):
\begin{equation}
\bigl\{ \xi \in \se(3) : \xi^\top (\Jpose \Jpose^\top)^{-1} \xi \leq 1 \bigr\},
\end{equation}
which characterizes end-effector twists reachable from a unit-norm joint velocity. Setting
\begin{equation}
W(q, z_e) = \bigl( \Jpose(q, z_e) \, \Jpose(q, z_e)^\top \bigr)^{-1}
\end{equation}
makes the metric $G_{\text{pose}}$ kinematic-dexterity-aware:
\begin{equation}
G_{\text{pose}}(q, z_e) = I_{n_q} + \Jpose^\top \bigl( \Jpose \Jpose^\top \bigr)^{-1} \Jpose.
\end{equation}
Near kinematic singularities ($\Jpose$ rank-deficient), $W$ diverges along the singular direction, so $G_{\text{pose}}$ assigns large metric distance to perturbations that approach the singularity. Forward diffusion under this $G_{\text{pose}}$ would naturally avoid singular configurations, and the volume correction $\tfrac{1}{2}\log\det G_{\text{pose}}$ in the limiting potential would induce a prior favoring manipulable regions of $\Mpose(z_e)$ \emph{without} an explicit singularity-avoidance term. The Riemannian guidance gradient
\begin{equation}
G_{\text{pose}}^{-1} \nabla_q R = \bigl( I + \Jpose^\top (\Jpose \Jpose^\top)^{-1} \Jpose \bigr)^{-1} \nabla_q R
\end{equation}
recovers a manipulability-weighted variant of the dynamically consistent damped pseudo-inverse familiar from operational-space control (Khatib, 1987; Nakamura, 1991).

\subsection{Choice (ii): joint-space inertia}

In rigid-body dynamics, the joint-space inertia matrix $M_\text{inert}(q) \in \R^{n_q \times n_q}$ defines the kinetic-energy quadratic form $T_{\text{kin}} = \tfrac{1}{2} \dot q^\top M_\text{inert}(q) \dot q$. Replacing the chart-coordinate identity $I_{n_q}$ in $G_{\text{pose}}$ with $M_\text{inert}(q)$ yields
\begin{equation}
G_{\text{pose}}^{\text{phys}}(q, z_e) = M_\text{inert}(q) + \Jpose(q, z_e)^\top W \, \Jpose(q, z_e),
\end{equation}
turning the Riemannian Gaussian noise $\xi^q \sim \N(0, (G_{\text{pose}}^{\text{phys}})^{-1})$ into an energy-aware perturbation: heavy joints receive smaller chart-space step variance, light joints receive larger, in proportion to inverse inertia. This connects diffusion sampling to the intrinsic Riemannian metric of configuration space induced by the kinetic energy, the standard geometric framework in classical mechanics (Arnold, 1989). It also relates to the literature on covariant Hamiltonian Monte Carlo (Girolami \& Calderhead, 2011), in which the local metric is chosen as the Fisher information or a related dynamical quantity.

\subsection{Choice (iii): task-space ellipsoid}

For tasks with anisotropic tolerance (e.g., peg-in-hole insertion: tight along the hole axis, loose laterally), $W_p$ can be chosen as a task-aligned ellipsoid:
\begin{equation}
W_p = R_{\text{task}} \, \mathrm{diag}\bigl( \sigma_x^{-2}, \sigma_y^{-2}, \sigma_z^{-2} \bigr) \, R_{\text{task}}^\top,
\end{equation}
with analogous structure for $W_R$. This is a strict generalization of the diagonal scalar choice and follows the pattern used in operational-space stiffness specification (Hogan, 1985). Within the present framework, $W_\text{task}$ would be a per-task input alongside $T_{\text{target}}$, modulating both the limiting distribution and the guidance gradient.

\subsection{Implications and open questions}

Each of the above choices preserves all structural properties established in this document: tangent verification on $\Mpose$, manifold adherence by construction, score lift via $J_H^{\text{pose}}$, DSM target via weighted pseudoinverse, and trajectory product manifold structure. The only consequence of upgrading $W$ from constant diagonal to $q$-dependent is that the volume correction $\tfrac{1}{2}\log\det G_{\text{pose}}(q, z_e)$ and the chart drift $b^q$ acquire additional terms reflecting $\partial_q W$. The autograd-based implementation of $\nabla_q R$ and $\nabla_q U$ accommodates this without code changes.

Open questions that future work should address:
\begin{enumerate}[leftmargin=*, noitemsep]
    \item Does manipulability-weighted $W$ improve performance on tasks where singularity avoidance matters (e.g., trajectories near a workspace boundary, dexterity-critical motions)?
    \item Does inertia-weighted $W$ produce smoother or more energy-efficient trajectories than chart-Euclidean $I_{n_q}$, particularly on under-actuated or compliance-rich systems?
    \item How does the $q$-dependence of $W$ affect the conditioning of $G_{\text{pose}}$, and what numerical stabilization (e.g., damped variants $W + \lambda I$) is needed in practice?
    \item Can the framework adopt task-conditional $W$ as part of the goal specification $c$, treating task ellipsoids as another modality of conditioning?
\end{enumerate}

These extensions are consistent with the framework's mathematical core (Sections~1--9) and instantiate the application-specific layer (Section~11) with richer task and dynamics structure. We leave their empirical investigation to future work.

\appendix

\section{Closed-form reward gradients (for reference; verify before use)}
\label{app:closed-form-grad}

This appendix documents closed-form expressions for the pose reward gradients $\nabla_{q_h} R_{\text{start/goal}}$. \textbf{The main implementation uses automatic differentiation rather than these closed forms}, due to (i) sensitivity of sign and Jacobian convention to the choice of left/right trivialization and the perturbation side, and (ii) the modest computational overhead of autograd through $\Log_{\SE(3)}$. The closed forms below are provided for cross-verification with autograd output and for understanding the gradient structure, but should be checked numerically before being relied upon directly.

\subsection{Differential of $\Log_{\SE(3)}$}

For $\xi \in \se(3)$, the $\SE(3)$ left Jacobian $\mathcal{J}_l(\xi) \in \R^{6 \times 6}$ has block form
\begin{equation}
\mathcal{J}_l(\xi) = \begin{bmatrix} J_{l, \SO(3)}(\omega) & Q(\rho, \omega) \\ 0 & J_{l, \SO(3)}(\omega) \end{bmatrix}, \qquad \xi = (\rho, \omega),
\end{equation}
with $J_{l, \SO(3)}(\omega) = I + \tfrac{1 - \cos \theta}{\theta^2}[\omega]_\times + \tfrac{\theta - \sin \theta}{\theta^3}[\omega]_\times^2$, $\theta = \|\omega\|$, and $Q(\rho, \omega)$ the $\SE(3)$ coupling block (Barfoot, \emph{State Estimation for Robotics}, Eq.~7.85b). The relevant linearization, depending on the perturbation side, is one of:
\begin{align}
\Log_{\SE(3)}\bigl( \exp(\xi^\wedge) \exp(\delta^\wedge) \bigr) &\approx \xi + \mathcal{J}_l^{-1}(\xi) \delta && \text{(right perturbation, } \mathcal{J}_l^{-1}\text{)} \\
\Log_{\SE(3)}\bigl( \exp(\delta^\wedge) \exp(\xi^\wedge) \bigr) &\approx \xi + \mathcal{J}_r^{-1}(\xi) \delta && \text{(left perturbation, } \mathcal{J}_r^{-1}\text{)}
\end{align}
with $\mathcal{J}_r(\xi) = \mathcal{J}_l(-\xi)$.

\subsection{Endpoint goal gradient — verified closed form}

\textbf{Resolved (numerical verification, Appendix~\ref{app:verify}, \texttt{tests/test\_reward\_gradient\_verification.py}):} sign $= +$, Jacobian $= \mathcal{J}_l$ (left). Relative error vs autograd: $9.67 \times 10^{-5}$ at $\|e_{\text{goal}}\|_\infty = 0.33$ rad on Franka 7-DoF.

For $e_{\text{goal}}(q_H) = \Log_{\SE(3)}\bigl( \Tphi(q_H, z_e)^{-1} T_{\text{target}} \bigr)$ with $W_g = \mathrm{diag}(\alpha_p^{\text{g}} I_3, \alpha_R^{\text{g}} I_3)$, the verified gradient is
\begin{equation}
\boxed{ \nabla_{q_H} R_{\text{goal}} \;=\; +2 \, \Jpose(q_H, z_e)^\top \, \mathcal{J}_l^{-\top}(e_{\text{goal}}) \, W_g \, e_{\text{goal}}(q_H, z_e). }
\end{equation}
\textbf{Derivation.} Body-frame perturbation $q_H \to q_H + \delta q$ right-multiplies $\Tphi$ by $\exp(\Jpose \delta q)^\wedge$, so $\Tphi^{-1} T_{\text{target}}$ is \emph{left}-multiplied by $\exp(-\Jpose \delta q)^\wedge$. Using the left-perturbation linearization $\Log_{\SE(3)}(\exp(\delta) \exp(\xi)) \approx \xi + \mathcal{J}_r^{-1}(\xi) \delta$ with $\delta = -\Jpose \delta q$, $\xi = e_{\text{goal}}$:
\begin{equation}
\partial_q e_{\text{goal}} = -\mathcal{J}_r^{-1}(e_{\text{goal}}) \, \Jpose.
\end{equation}
Differentiating $R_{\text{goal}} = -e_{\text{goal}}^\top W_g e_{\text{goal}}$:
\begin{equation}
\nabla_{q_H} R_{\text{goal}} = -2 e_{\text{goal}}^\top W_g \cdot (-\mathcal{J}_r^{-1}) \Jpose = +2 \Jpose^\top \mathcal{J}_r^{-\top}(e_{\text{goal}}) W_g e_{\text{goal}}.
\end{equation}
The verification protocol identifies the matching $\mathcal{J}_\bullet$ as $\mathcal{J}_l$, not $\mathcal{J}_r$, indicating that our $\mathcal{J}$ block construction (Barfoot Eq.~7.85b) corresponds to the left-Jacobian under the framework's sign convention; the numerical match confirms $+\mathcal{J}_l$ at $9.67\times10^{-5}$ rel.\ error vs autograd.

\subsection{Decoupled form (recommended for paper text)}

If the reward uses the decoupled form $R_{\text{goal}} = -\alpha_p \| p_\phi(q_H) - p_{\text{target}} \|^2 - \alpha_R \| \Log_{\SO(3)}(R_\phi(q_H)^\top R_{\text{target}}) \|^2$, the position term is purely Euclidean and unambiguous:
\begin{equation}
\nabla_{q_H} R_{\text{goal}}^{\text{(pos part)}} = -2 \alpha_p \, J_p^{\text{world}}(q_H, z_e)^\top \, (p_\phi(q_H, z_e) - p_{\text{target}}),
\end{equation}
where $J_p^{\text{world}} = \partial p_\phi / \partial q$ is the world-frame position Jacobian. The rotation term involves the $\SO(3)$ left Jacobian:
\begin{equation}
\nabla_{q_H} R_{\text{goal}}^{\text{(rot part)}} \;\overset{?}{=}\; \pm 2 \alpha_R \, J_\omega(q_H, z_e)^\top \, J_{l, \SO(3)}^{-\top}(e_\omega^g) \, e_\omega^g, \qquad e_\omega^g = \Log_{\SO(3)}(R_\phi(q_H)^\top R_{\text{target}}).
\end{equation}
\verifyflag{Sign of the rotation gradient depends on perturbation side; verify numerically.}

\subsection{PyTorch verification protocol}
\label{app:verify}

To resolve all $\pm$ and $\mathcal{J}_l$/$\mathcal{J}_r$ ambiguities, run the following protocol on a small example before deploying any closed-form gradient. The autograd output is taken as ground truth.

\begin{enumerate}[leftmargin=*]
    \item \textbf{Setup}: Pick random $q_H \in \R^{n_q}$, random $T_{\text{target}} \in \SE(3)$, fixed $z_e$. Set $\alpha_p^g, \alpha_R^g$.
    \item \textbf{Autograd reference}: Compute $R_{\text{goal}}(q_H)$ as a scalar via PyTorch operations on $\Tphi$ and the chosen log map, and obtain $g_{\text{ag}} = \nabla_{q_H} R_{\text{goal}}$ via \texttt{torch.autograd.grad}.
    \item \textbf{Finite-difference reference}: For each coordinate $i$, compute
    $g_{\text{fd}, i} \approx [R_{\text{goal}}(q_H + \epsilon e_i) - R_{\text{goal}}(q_H - \epsilon e_i)] / (2 \epsilon)$
    with $\epsilon = 10^{-5}$. Confirm $\| g_{\text{ag}} - g_{\text{fd}} \|_\infty < 10^{-4}$.
    \item \textbf{Closed-form candidate}: Compute the candidate closed-form $g_{\text{cf}}$ for each combination $(\text{sign}, \mathcal{J}_\bullet) \in \{+, -\} \times \{\mathcal{J}_l, \mathcal{J}_r\}$. Identify the combination matching $g_{\text{ag}}$ (i.e., $\| g_{\text{cf}} - g_{\text{ag}} \|_\infty < 10^{-4}$).
    \item \textbf{Approximation regime}: For the simplified gradient $g_{\text{simple}} = -2 \Jpose^\top W_g e_{\text{goal}}$ (no $\mathcal{J}^{-\top}$ factor, sign as in v1 of this document), measure $\| g_{\text{simple}} - g_{\text{ag}} \|$ as a function of $\| e_{\text{goal}} \|$. Determine the threshold $\| e_{\text{goal}} \|^*$ above which the simplified form deviates by more than (e.g.) 5\%.
    \item \textbf{Document outcome}: Update Section~\ref{app:closed-form-grad} of this document with the verified sign and Jacobian convention. Also update the small-error approximation guidance with the measured threshold.
\end{enumerate}

\textbf{Recommended action}: Until step 6 is completed, use autograd as the production gradient implementation. The closed forms above are documentation, not code.

\subsection{Self-consistency check on simple cases}

Two limit cases provide additional cross-checks:

\paragraph{(i) Pure position task ($\alpha_R^g = 0$, $T_{\text{target}}$ has $R_{\text{target}} = R_\phi(q_H)$).} The body-frame translation error reduces to $e_\rho^g = R_\phi^\top (p_\text{target} - p_\phi)$, and the gradient should match the position-only framework's $-2 \alpha_p^g \Jpose^\top W_g (\cdot)$ structure (after accounting for the body-frame $R^\top$ rotation). This recovers the v1 position-only formulation.

\paragraph{(ii) Vanishing error ($e_{\text{goal}} \to 0$).} The gradient approaches zero, and the $\mathcal{J}^{-\top}(e)$ factor approaches $I_6$. The simplified form $+2 \Jpose^\top W_g e$ (\textbf{verified $+$ sign}) is exact in this limit. Empirical $5\%$ threshold (\texttt{test\_reward\_gradient\_verification.py} sweep on Franka): $\|e\|_\infty \lesssim 0.75$ rad (rotation-dominated). Beyond this regime, the $\mathcal{J}_l^{-\top}$ factor must be retained.

\section{Summary of fixes from v1 to v2}

\begin{center}
\begin{tabular}{l|p{4cm}|p{4cm}}
\textbf{Fix} & \textbf{v1} & \textbf{v2} \\
\hline
Metric form & $G_{\text{pose}} = (J_H^{\text{pose}})^\top W J_H^{\text{pose}}$ (dimension mismatch: $W$ is $6\times 6$, $J_H^{\text{pose}}$ is $(n_q+6) \times n_q$) & $G_{\text{pose}} = I_{n_q} + \Jpose^\top W \Jpose$ (dimensionally consistent: $\Jpose$ is $6 \times n_q$) \\
\hline
DSM target weight & $W$ multiplied with $(n_q + 6)$-vector $\Exp^{-1}$ (dimension mismatch) & Block decomposition $a^* = G^{-1}[(q_0-q_r) + \Jpose^\top W \Log_{\SE(3)}(\cdot)]$, equivalent to $\widetilde W = \mathrm{diag}(I_{n_q}, W)$ \\
\hline
Log map naming & $\Exp^{-1}_{x_r}(x_0)$ written as if intrinsic Riemannian log & $\delta x^{\text{amb}}_{r0}$ explicitly named as ambient/trivialized displacement; intrinsic log denoted $\mathrm{Log}_{\Mset, x_r}$; approximation made explicit \\
\hline
Chart target identity & $J_H^{\text{pose}} a^* = \Exp^{-1}_{x_r}(x_0)$ written as equality & $J_H^{\text{pose}} a^* = \Pi^{\widetilde W}(\delta x^{\text{amb}})$ as projection (equality does not hold in general) \\
\hline
Reward gradient & Closed-form with sign $-2 \Jpose^\top \mathcal{J}_l^{-\top} W e$ in main text & Autograd in main text; closed forms moved to Appendix~\ref{app:closed-form-grad} with \textbf{[VERIFY]} flags \\
\hline
``Ambient dimension'' & $n_q + 6$ as ``ambient dimension'' & ``trivialized tangent / configuration manifold dimension'' (distinct from $\SE(3)$'s matrix-embedding dimensions) \\
\hline
Position-only specialization & ``exactly recovers'' position-only & Up to choice of residual frame; $\Delta_\phi = R_{\text{analytic}} \rho_\phi$ relation made explicit \\
\end{tabular}
\end{center}

\paragraph{Verification status.} The mathematical structure of v2 is dimensionally consistent and conceptually well-formed. Verification items:
\begin{enumerate}[leftmargin=*, noitemsep]
    \item \textbf{[RESOLVED]} Reward gradient sign and $\mathcal{J}_l$/$\mathcal{J}_r$ convention — sign $=+$, Jacobian $=\mathcal{J}_l$ (\texttt{tests/test\_reward\_gradient\_verification.py}, rel.\ error $9.67\times10^{-5}$ vs autograd; simplified form $5\%$ threshold $\|e\|\lesssim0.75$ rad).
    \item Pose self-model fitting accuracy on synthetic compliance — implemented (\texttt{smcdp/experiments/franka\_stage1\_selfmodel\_pose.py}); empirical numbers pending GPU run.
    \item Pose-target task experiments — implemented (\texttt{smcdp/experiments/franka\_traj\_unet\_pose.py}); empirical numbers pending GPU run.
\end{enumerate}

\end{document}